% SPDX-License-Identifier: GPL-3.0-or-later
% Copyright (C) 2026 Alessandro Greco
\chapter{Ошибки операционной безопасности и деанонимизация: уроки}\label{cap:28-opsec}
% Колонтитул короче: полное название почти касается номера страницы.
\chaptermark{Ошибки операционной безопасности и деанонимизация}
\aggiornatoal{23 сентября 2026 г.}{рекомендации Tor Project; предупреждения Tails; tor 0.4.9.13; лабораторная сеть}

\begin{obiettivi}
К концу этой главы вы сможете:
\begin{itemize}[nosep]
  \item различать уровни, на которых может произойти утечка идентичности, и определять, какие инструменты защищают каждый из них;
  \item распознавать распространённые операционные ошибки — от утечек DNS до слишком малого множества анонимности;
  \item использовать сам tor, чтобы проверить, передаёт ли приложение ему разрешение имён (Лабораторная работа~18);
  \item анализировать реальный инцидент в терминах модели угроз.
\end{itemize}
\end{obiettivi}

\Cref{cap:26-attacchi} завершалась неприятным наблюдением: атаки на протокол могут быть дорогими и требовать привилегированных точек наблюдения. Во многих описанных случаях утрата анонимности также связана с ошибками использования инструментов, которые один лишь сетевой протокол предотвратить не может. В этой главе рассматриваются такие ошибки — не для того, чтобы объяснять, как идентифицировать людей, а чтобы ясно провести границу между защитой Tor и ответственностью самого пользователя.

\begin{attenzione}[Границы этой главы]
Мы обсуждаем ошибки и защитные практики с точки зрения людей, стремящихся защитить себя. Процедуры идентификации людей не приводятся. Упомянутые инциденты являются публичными и реконструированы по указанным источникам.
\end{attenzione}

% ---------------------------------------------------------------------------
\section{Четыре уровня}\label{sec:opsec-livelli}

Идентичность может утечь на четырёх уровнях (\cref{fig:livelli}). Первый — сеть — это уровень, который защищает tor: при обычном соединении через Tor назначение не видит исходный IP-адрес клиента, а отдельный ретранслятор цепи обычно не обладает полной картиной обоих концов. Это не исключает корреляционные атаки и другие риски. Второй — приложение — рассматривается Tor Browser, Tails и Whonix (\cref{cap:23-browser}). Последние два — содержимое и поведение — tor не защищает автоматически. Специальные инструменты способны снизить часть рисков, но важно и то, что пользователь пишет, загружает и делает.

\begin{figure}[htbp]
\centering
% SPDX-License-Identifier: GPL-3.0-or-later
% Copyright (C) 2026 Alessandro Greco
%% Русское альтернативное описание EPUB.
\altfig{Четыре уровня возможного раскрытия личности: сеть (IP-адреса и цепи), приложение (DNS, плагины и цифровые отпечатки), содержимое (метаданные, учётные записи и стиль), поведение (время, привычки и контакты). Tor снижает некоторые сетевые риски; специальные системы и инструменты помогают на других уровнях наряду с решениями пользователя. Ни один уровень сам по себе не гарантирует анонимности.}
% Уровни утечек идентичности и соответствующие защиты
\begin{tikzpicture}[font=\sffamily\scriptsize,
  liv/.style={draw,rounded corners=2pt,minimum width=5.2cm,minimum height=0.75cm,align=center},
  chi/.style={align=left,text width=5cm}]
  \node[liv,fill=torviola!15] (r) at (0,0) {сеть: IP-адреса и цепи};
  \node[liv,fill=torviola!9] (a) at (0,0.95) {приложение: DNS, плагины, отпечатки};
  \node[liv,fill=torgrigio!10] (c) at (0,1.9) {содержимое: метаданные, аккаунты, стиль};
  \node[liv,fill=avviso!12] (b) at (0,2.85) {поведение: время, привычки, контакты};
  \node[chi,right=0.4cm of r] {tor: сетевая защита не абсолютна};
  \node[chi,right=0.4cm of a] {Tor Browser, Tails, Whonix (\cref{cap:23-browser})};
  \node[chi,right=0.4cm of c] {инструменты и выбор содержимого};
  \node[chi,right=0.4cm of b] {решения, привычки, другие инструменты};
  \draw[->,thick,torgrigio] (-3.1,-0.3) -- node[left,rotate=90,anchor=south] {защита не полностью автоматическая} (-3.1,3.2);
\end{tikzpicture}

\caption{Учебная схема четырёх уровней возможного раскрытия личности. Tor снижает часть сетевых рисков, приложения и инструменты обработки содержимого помогают, но раскрываемые сведения и поведение важны на каждом уровне. Ни один уровень сам по себе не гарантирует анонимности.}
\label{fig:livelli}
\end{figure}

% ---------------------------------------------------------------------------
\section{Ошибки приложений}\label{sec:opsec-applicazione}

\paragraph{Трафик вне Tor.} tor защищает только приложения, настроенные на его использование~\cite{torsafety}. Программа, выполняющая прямые соединения, даже ради обновлений или предпросмотра, может раскрыть реальный IP-адрес. Служба поддержки Tor Project выделяет два известных случая. BitTorrent-приложения могут игнорировать настройки прокси и дополнительно раскрывать адрес трекерам. Загруженные документы, открытые при активном сетевом соединении, создают другой риск: DOC- или PDF-файлы могут содержать удалённые ресурсы, которые загружаются приложениями, не использующими Tor.

\paragraph{Утечки DNS.} Более тонкий случай связан с SOCKS. Приложение может передать tor \emph{имя} назначения, позволяя выходному ретранслятору разрешить его (\cref{sec:stream-dns}), либо самостоятельно разрешить имя через системную DNS-службу и передать tor только IP-адрес. Во втором случае последующее соединение всё ещё идёт через Tor, но доменное имя могло уже раскрыться системе разрешения имён вне цепи Tor и оказаться связано с сетевым подключением клиента. Открытый DNS-запрос не обязателен: адрес мог быть получен из кеша, локального файла или по зашифрованному DNS. tor не может доказать, что утечка уже произошла, но способен предупредить о SOCKS-запросе только с IP-адресом, как показывает лабораторная работа.

% ---------------------------------------------------------------------------
\section{Лабораторная работа: tor обнаруживает признак DNS leak}\label{sec:lab-opsec}

\begin{laboratorio}[18]{SOCKS с hostname и без него}
\textbf{Цель.} Наблюдать, как tor различает приложение, передающее ему разрешение имени, и приложение, выполняющее разрешение самостоятельно, а также как tor может отклонить второй вариант.\\
\textbf{Предварительные условия.} Лабораторная работа~11; лабораторная сеть должна быть запущена.\\
\textbf{Оценочное время.} 30 минут.
\end{laboratorio}

\cmd{curl} поддерживает оба варианта. При \texttt{-{}-socks5} он разрешает имя локально и передаёт tor IP-адрес; при \texttt{-{}-socks5-hostname} передаёт само имя. Скрипт \file{labs/ru/лаборатория/dns\_leak.py} подключается к ControlPort клиента \texttt{008c} и включает \opt{TestSocks}, заставляя tor записывать тип каждого SOCKS-запроса. Скрипт слушает сообщения журнала и проверяет оба варианта на \texttt{http://localhost:8080/}.

Эта лабораторная работа демонстрирует \emph{механизм}, а не реальную DNS leak, наблюдаемую на линии. \texttt{localhost} разрешается локально без обращения к resolver, поэтому запрос не покидает компьютер. tor получает IP-адрес вместо имени хоста: это возможный признак предварительного локального разрешения, а не доказательство утечки DNS (программа могла сразу использовать IP-адрес). Если использовать обычный домен вместо \texttt{localhost}, аналогичное локальное разрешение может породить внешний DNS-запрос, который при определённой конфигурации будет открытым. Мы намеренно не генерируем такой трафик и остаёмся внутри приватной сети. Скрипт также принудительно использует SOCKS для \texttt{localhost} вопреки возможному исключению \texttt{NO\_PROXY}. Наконец, он включает \opt{SafeSocks}, отклоняющий запросы только с IP-адресом даже при отсутствии DNS-утечки:
\begin{lstlisting}[numbers=none]
--socks5           HTTP 200  tor: Your application (using socks5 to port
                   8080) is giving Tor only an IP address. Applications
                   that do DNS resolves themselves may leak information...
--socks5-hostname  HTTP 200  tor: Your application (using socks5 to port
                   8080) instructed Tor to take care of the DNS resolution
                   itself if necessary. This is good.
--socks5 SafeSocks HTTP 000
\end{lstlisting}
В иллюстративном выводе оба первых запроса успешны, но только второй передаёт разрешение имени tor; первый не доказывает DNS-утечку, поскольку используется \texttt{localhost}. В коде tor (\file{src/core/proto/proto\_socks.c}) первое предупреждение записывается на уровне \texttt{warn}, с ограничением частоты записи, чтобы не переполнять журнал, и сопровождается событием \texttt{STATUS\_CLIENT} с действием \texttt{DANGEROUS\_SOCKS} на ControlPort. Поэтому управляющее приложение может предупредить пользователя. При включённом \opt{SafeSocks} тот же запрос отклоняется, и \cmd{curl} не получает ответа. Лабораторная работа включена в \texttt{make verifica-lab}.

% ---------------------------------------------------------------------------
\section{Ошибки содержимого}\label{sec:opsec-contenuto}

\paragraph{Метаданные.} Многие файлы содержат информацию, которую авторы даже не замечают: фотографии могут раскрывать местоположение или модель камеры; документы — автора и даты создания. Tails включает инструмент удаления метаданных, а его предупреждения упоминают случаи, когда метаданные фотографии раскрывали, где находился фотограф~\cite{tailswarnings}.

\paragraph{Учётные записи и идентичность.} Вход в личную учётную запись сообщает сайту, кто вы, независимо от видимого IP-адреса соединения~\cite{torsafety}. Использование двух идентичностей на одном сайте в одном сеансе также может связать их через cookie, данные сеанса, идентификаторы приложений и другие признаки. При этом сайт не видит напрямую идентификатор цепи Tor. Tails рекомендует перезапускаться между действиями, которые должны оставаться разделёнными~\cite{tailswarnings}. Требование несвязываемости Tor Browser (\cref{sec:browser-tb}) относится к разным сайтам, а не к разным идентичностям на одном сайте.

\paragraph{Стиль письма.} Даже без метаданных особенности письма, тематика, повторяющиеся ошибки и время публикаций могут связать анонимный текст с подписанным. Ни один автоматический инструмент не способен устранить все такие признаки.

% ---------------------------------------------------------------------------
\section{Поведенческие ошибки: слишком малое множество анонимности}\label{sec:opsec-insieme}

В \cref{cap:01-anonimato} анонимность определялась через неразличимость внутри множества. Tor может уменьшить некоторые признаки, позволяющие различать сетевую активность, но не гарантирует, что в конкретном месте и в конкретное время Tor использует достаточно людей. Если пользователей мало, знание того, \emph{кто использовал Tor}, может существенно сузить множество вариантов.

\paragraph{Инцидент.} В декабре 2013 года один американский университет получил ложные угрозы взрыва по электронной почте незадолго до экзамена. Отправитель использовал временный почтовый сервис и Tor. Согласно письменному показанию под присягой сотрудника ФБР, приложенному к обвинению и описанному несколькими СМИ, университет определил, кто использовал Tor в его беспроводной сети в часы перед отправкой сообщений~\cite{wbur2013harvard,sandvik2013harvard}. Студент позднее признался во время допроса. Для сужения круга подозреваемых не понадобилось взламывать шифрование Tor: локальные наблюдатели могли установить факт использования Tor, что само по себе не означало знания содержимого или назначения соединений. Но они могли видеть, \emph{кто} использовал Tor в этот временной интервал, уменьшив круг подозреваемых до нескольких имён.

Инцидент иллюстрирует три принципа модели угроз (\cref{cap:04-modello-minaccia}):
\begin{itemize}
  \item Tor без моста не скрывает сам факт использования: наблюдатели локальной сети могут распознать его~\cite{tailswarnings}. Мост с подключаемым транспортом (\cref{cap:16-censura}) способны уменьшить эту заметность.
  \item Корреляции не обязательно нужен глобальный противник, если множество кандидатов мало, а событие происходит в узком временном окне. Это базовая задача пересечения множеств.
  \item Анонимность любит компанию~\cite{dingledine2006company}. Чем больше людей с разными причинами используют Tor в одной среде, тем меньше информации несёт сам факт использования конкретным человеком. Служба поддержки Tor Project прямо поощряет более широкое распространение Tor~\cite{torsafety}.
\end{itemize}

% ---------------------------------------------------------------------------
\section{Контрольный список}\label{sec:opsec-lista}

Список не исчерпывающий, но последовательно охватывает уровни:
\begin{enumerate}
  \item Используйте Tor Browser или систему вроде Tails или Whonix, а не обычный браузер с вручную настроенным прокси.
  \item Не устанавливайте расширения и не меняйте настройки так, чтобы браузер выделялся.
  \item Для других приложений проверяйте, что они передают имена хостов в tor (\opt{SafeSocks}, \opt{TestSocks}).
  \item Не открывайте загруженные документы при подключении к сети либо используйте для них изолированную среду.
  \item Предпочитайте HTTPS и onion-сервисы, поскольку выходной ретранслятор может видеть незашифрованное содержимое.
  \item Удаляйте метаданные из файлов перед публикацией.
  \item Не объединяйте в одном сеансе идентичности, которые должны оставаться разделёнными.
  \item Если сам факт использования Tor создаёт риск, рассмотрите подходящий замаскированный мост, не предполагая полной невидимости.
  \item Помните, что расписание, привычки и стиль письма тоже передают информацию.
\end{enumerate}

% ---------------------------------------------------------------------------
\section*{Итоги}
\begin{itemize}
  \item Идентичность может утечь на сетевом, прикладном, содержательном и поведенческом уровнях. tor напрямую защищает только первый.
  \item Приложения, обходящие Tor, загруженные документы, открытые онлайн, и DNS leaks — ошибки уровня приложения. tor может предупреждать о SOCKS-запросах, содержащих только IP-адрес, и отклонять их.
  \item Метаданные, учётные записи и стиль письма способны связать содержимое с человеком.
  \item Когда множество пользователей Tor мало, знания о том, кто и когда использовал Tor, может быть достаточно для существенного сужения кандидатов. Мосты и большое сообщество пользователей помогают снизить риск.
\end{itemize}

% ---------------------------------------------------------------------------
\section*{Упражнения}

\begin{esercizio}[1]\label{es:28-livello}
Для каждой ошибки, рассмотренной здесь, определите её уровень на \cref{fig:livelli} и инструмент либо привычку, помогающие её предотвратить.
\end{esercizio}

\begin{esercizio}[1]\label{es:28-dns}
Объясните, почему DNS leak раскрывает информацию, даже если последующее соединение проходит через Tor. Кто видит эту информацию?
\end{esercizio}

\begin{esercizio}[2]\label{es:28-evento}
В Лабораторной работе~18 подпишитесь через Stem на события \texttt{STATUS\_CLIENT} и выведите действие \texttt{DANGEROUS\_SOCKS}, создаваемое запросом только с IP. Как графический интерфейс мог бы использовать это событие?
\end{esercizio}

\begin{esercizio}[2]\label{es:28-insieme}
Повторно рассмотрите инцидент 2013 года, предположив, что студент использовал bridge obfs4 из той же сети. Что могли бы видеть наблюдатели локальной сети и какая другая информация всё ещё могла сузить множество кандидатов?
\end{esercizio}

\begin{esercizio}[3]\label{es:28-modello}
Выберите реального или вымышленного человека с законной потребностью в анонимности, например журналистский источник, активиста или исследователя, и составьте операционный план: угрозы, соответствующие уровни, инструменты и привычки, остаточные риски.
\end{esercizio}

% ---------------------------------------------------------------------------
\section*{Дополнительная литература}
\begin{itemize}
  \item Рекомендации службы поддержки Tor Project по анонимности~\cite{torsafety}.
  \item Предупреждения Tails, \titolo{Tails is safe but not magic}~\cite{tailswarnings}.
  \item R.~Dingledine и N.~Mathewson, \titolo{Anonymity Loves Company}~\cite{dingledine2006company}.
\end{itemize}
